\section{ویژگی‌های ذره‌ها}
\label{sec:particles}
در این بخش هر ۱۷ ویژگی ذرات تشکیل‌دهنده (تعریف‌ها در بخش~\ref{sec:vars}) برای ۱۰ ذره پرانرژی‌تر هر جت بررسی می‌شود. داده، همان داده نهایی (کامل و جزئی) است. هر شکل ده نمودار دارد، یکی برای هر ذره به ترتیب انرژی، و در زیر هر نمودار نسبت دو جفت نمونه رسم شده است: نسبت \Zjet به \bbbb (بنفش) و نسبت \ttZ به \ttbb (نارنجی). مقدارهای عددی این بخش را از روی محور شکل‌ها می‌خوانیم و با «حدود» می‌آوریم، چون خطای آماری هر بین محاسبه نشده است و نسبت‌های نزدیک به ۱ در دنباله توزیع‌ها نوسان آماری دارند. ویژگی‌ها را در چهار گروه می‌آوریم: تکانه و انرژی در دستگاه آزمایشگاه، کمیت‌های نسبی، زاویه‌ها و نوع ذره. در هر گروه به بخش‌های قبلی که با آن مرتبط‌اند ارجاع می‌دهیم.

\subsection{تکانه و انرژی در دستگاه آزمایشگاه}
\widefig{cons_01.pdf}{$p_x$ ده ذره پرانرژی.}{fig:cons-px}
\widefig{cons_02.pdf}{$p_y$ ده ذره پرانرژی.}{fig:cons-py}
توزیع $p_x$ و $p_y$ (شکل‌های~\ref{fig:cons-px} و~\ref{fig:cons-py}) دو قله در دو طرف صفر دارد و عرض آن‌ها با رتبه ذره کم می‌شود: دامنه محور برای ذره اول حدود $\pm150$ و برای ذره دهم حدود $\pm15$ است. این شکل با یک استدلال ساده توضیح داده می‌شود. اگر ذره تقریباً در راستای جت حرکت کند، آنگاه $p_x\approx p_{T,i}\cos\varphi$ است که $\varphi$ زاویه سمتی جت است. زاویه سمتی جت‌ها یکنواخت است (شکل~\ref{fig:cons-phi} در ادامه)، و چگالی $x=\cos\varphi$ برای $\varphi$ یکنواخت برابر $1/(\pi\sqrt{1-x^2})$ است که در $x=\pm1$ بی‌نهایت می‌شود. پس دو قله در $p_x\approx\pm p_{T,i}$ انتظار می‌رود و پهن‌شدن آن‌ها از پراکندگی $p_{T,i}$ می‌آید. $p_y$ با $\sin\varphi$ همین شکل را دارد. نسبت‌ها در بخش مرکزی حدود ۱ هستند و فقط در دنباله، به‌ویژه برای ذره اول، نسبت \Zjet به \bbbb تا حدود ۲ بالا می‌رود. این بالا رفتن با افزایش نسبت در $p_T$ (شکل~\ref{fig:cons-pt}) سازگار است. چون $p_x$ و $p_y$ فقط جهت سمتی جت را به $p_T$ ذره اضافه می‌کنند و جهت سمتی هیچ اطلاعاتی ندارد، اطلاعات تازه‌ای نسبت به $p_T$ در این دو ویژگی دیده نمی‌شود.

\widefig{cons_03.pdf}{$p_z$ ده ذره پرانرژی.}{fig:cons-pz}
\widefig{cons_04.pdf}{$E$ ده ذره پرانرژی.}{fig:cons-e}
شکل~\ref{fig:cons-pz} متفاوت است. چون $p_z=p_T\sinh\eta$ است، توزیع $p_z$ به $\eta$ حساس است. نسبت \Zjet به \bbbb از حدود ۱ در مرکز به حدود ۲ در دنباله می‌رسد و برعکس، نسبت \ttZ به \ttbb در دنباله به حدود \lr{0.7} کاهش می‌یابد. پس دو جفت در خلاف جهت هم‌اند: سیگنال \Zjet پهن‌تر از \bbbb است ولی سیگنال \ttZ باریک‌تر از \ttbb. همین الگو در $\eta$ جت‌ها در بخش~\ref{sec:cutflow} و در بخش~\ref{sec:set4} (خطر میان‌بر) آمده و منشأ احتمالی آن در بخش~\ref{sec:zorigin} بحث شده است. انرژی (شکل~\ref{fig:cons-e}) رابطه $E=p_T\cosh\eta$ دارد و ترکیبی از اثر $p_T$ و $\eta$ است: نسبت \Zjet به \bbbb از ۱ در انرژی‌های کم تا حدود ۲ تا \lr{2.5} در انرژی‌های بالا می‌رود و نسبت \ttZ به \ttbb از ۱ تا حدود \lr{0.7} در انرژی بالا می‌افتد. چون جهت اثر در دو جفت مخالف است، این ویژگی‌ها نشانه \Zbos نیستند بلکه بازتاب تفاوت توزیع $\eta$ در نمونه‌ها هستند.

\widefig{cons_06.pdf}{$p_T$ ده ذره پرانرژی.}{fig:cons-pt}
\widefig{cons_07.pdf}{$p_T/p_{T,\mathrm{jet}}$ ده ذره پرانرژی.}{fig:cons-ptrel}
$p_T$ ذره (شکل~\ref{fig:cons-pt}) در نمونه‌ها تقریباً یکی است، با نسبتی که در قله حدود ۱ است. در دنباله، نسبت \Zjet به \bbbb تا حدود ۲ بالا می‌رود (برای ذره اول تا حدود \lr{2.5}) و نسبت \ttZ به \ttbb در بازه حدود ۱ تا \lr{1.5} و با نوسان زیاد می‌ماند. بخشی از این تفاوت می‌تواند از اختلاف طیف $p_T$ خود جت‌ها در بخش~\ref{sec:cutflow} بیاید. اما شکل~\ref{fig:cons-ptrel} که $p_T$ را بر $p_{T,\mathrm{jet}}$ تقسیم می‌کند و مقیاس جت را حذف می‌کند، همان بالا رفتن را نشان می‌دهد (نسبت از حدود \lr{0.8} در سمت چپ تا حدود ۲ در سمت راست برای ذره اول، و برای ذرات ۶ تا ۱۰ از حدود \lr{0.7} تا \lr{1.5}). \textcolor{red}{این بالا رفتن شایان توجه است و بهتر است دلیلش بررسی شود.}

\subsection{کمیت‌های نسبی}
\widefig{cons_05.pdf}{$E/E_\mathrm{jet}$ ده ذره پرانرژی.}{fig:cons-erel}
شکل~\ref{fig:cons-erel} تقریباً همان شکل~\ref{fig:cons-ptrel} است: هم توزیع‌ها و هم نسبت‌ها (از حدود \lr{0.5} تا \lr{0.8} در مقدارهای کم تا حدود \lr{1.5} تا ۲ در مقدارهای زیاد). این طبیعی است چون جت‌ها $|\eta|<2$ دارند و $E\approx p_T\cosh\eta$ برای ذرات داخل یک جت تقریباً با ضریب مشترک $\cosh\eta_\mathrm{jet}$ به $p_T$ مربوط است. پس $E/E_\mathrm{jet}$ اطلاعات مستقلی از $p_T/p_{T,\mathrm{jet}}$ نمی‌دهد. دو جفت در این ویژگی هم‌جهت‌اند: سهم ذره‌های سخت در سیگنال کمی بیشتر و سهم ذره‌های نرم کمتر است. \textcolor{red}{علت این تفاوت را نمی‌دانیم.}

\subsection{زاویه‌ها}
\widefig{cons_08.pdf}{$\eta$ ده ذره پرانرژی.}{fig:cons-eta}
\widefig{cons_11.pdf}{$\varphi$ ده ذره پرانرژی.}{fig:cons-phi}
$\eta$ ذره‌ها (شکل~\ref{fig:cons-eta}) توزیع $\eta$ جت‌ها را تکرار می‌کند، چون ذره‌ها در جت می‌مانند. نسبت \Zjet به \bbbb از حدود \lr{0.8} تا \lr{0.9} در مرکز به حدود \lr{1.3} تا \lr{1.6} در $|\eta|\approx2$ می‌رود، و نسبت \ttZ به \ttbb از حدود \lr{1.1} در مرکز به \lr{0.5} تا \lr{0.7} در لبه می‌افتد. این دو جفت مثل $p_z$ در خلاف جهت‌اند، پس باز با تفاوت توزیع $\eta$ جت‌ها (بخش‌های~\ref{sec:cutflow} و~\ref{sec:set4}) سازگارند. شکل~\ref{fig:cons-phi} نشان می‌دهد $\varphi$ در همه نمونه‌ها یکنواخت است و نسبت‌ها حدود ۱ می‌مانند، که انتظار می‌رود چون برخورد پروتون‌ها حول محور پرتو تقارن دارد. اولین و آخرین بین کمی بالاتر است (حدود ۱۰ تا ۲۰ درصد) که احتمالاً اثر لبه بین‌بندی است؛ این بررسی نشده است. این ویژگی را فقط به‌عنوان یک آزمون سلامت داده نگه می‌داریم.

\widefig{cons_09.pdf}{$\eta-\eta_\mathrm{jet}$ ده ذره پرانرژی.}{fig:cons-etarel}
\widefig{cons_12.pdf}{$\varphi-\varphi_\mathrm{jet}$ ده ذره پرانرژی.}{fig:cons-phirel}
\widefig{cons_14.pdf}{$\Delta R$ ذره تا محور جت.}{fig:cons-dr}
اگر محور جت را مبدأ بگیریم، سه شکل~\ref{fig:cons-etarel}، \ref{fig:cons-phirel} و~\ref{fig:cons-dr} با هم یک ماجرا را نشان می‌دهند، چون $\Delta R^2=(\eta-\eta_\mathrm{jet})^2+(\varphi-\varphi_\mathrm{jet})^2$. برای $\eta-\eta_\mathrm{jet}$ و $\varphi-\varphi_\mathrm{jet}$ هر دو قله مرکزی در پس‌زمینه بلندتر است (برای ذره اول حدود \lr{7.5} تا ۸ در برابر حدود \lr{4.5} در سیگنال) و نسبت سیگنال به پس‌زمینه دو برآمدگی در حدود $\pm0.1$ تا $\pm0.2$ با ارتفاع حدود \lr{1.5} تا \lr{1.8} و یک فرورفتگی در صفر (حدود \lr{0.7} تا \lr{0.8}) دارد. دو متغیر شکل یکسان دارند، که با تقارن دورانی پروفایل شعاعی ذرات حول محور جت سازگار است. در $\Delta R$ همین را به شکل دیگری می‌بینیم: قله $\Delta R$ در پس‌زمینه حدود \lr{0.05} و در سیگنال حدود \lr{0.15} تا \lr{0.2} است، نسبت سیگنال به پس‌زمینه در $0.1<\Delta R<0.25$ به حدود \lr{1.5} تا ۲ می‌رسد و برای $\Delta R<0.05$ و $\Delta R>0.35$ زیر ۱ است. دو جفت هم‌جهت‌اند.

این الگو با ساختار دو‌هسته‌ای سیگنال سازگار است. اگر $z$ سهم $p_T$ یکی از دو کوارک $b$ باشد، فاصله زاویه‌ای دو کوارک (بخش~\ref{sec:zorigin}) $\Delta R_{b\bar b}\approx m/(p_T\sqrt{z(1-z)})$ است و کمینه آن، در $z=1/2$، برابر $2m/p_T$ است؛ برای $m=91$ و $p_T=400$ این $2\times91/400\approx0.46$ می‌شود. برای آزمودن این استدلال، از ذره‌های پارتونی فایل‌های \lr{Delphes} (چهار کار از هر نمونه) $\Delta R_{b\bar b}$ را برای جت‌هایی که هر دو کوارک $b$ در آن‌ها است خواندیم؛ برای سیگنال این دو کوارک فرزندان $Z$ و برای پس‌زمینه دو کوارک $b$ پرانرژی‌تر داخل جت‌اند. میانه $\Delta R_{b\bar b}$ برای \Zjet و \ttZ برابر \lr{0.42} و \lr{0.40} و برای \bbbb و \ttbb برابر \lr{0.23} و \lr{0.24} است. میانه سیگنال با کمینه $2m/p_T\approx2\times91/490\approx0.37$ (برای $p_T$ میانه حدود ۴۹۰) سازگار است. در پس‌زمینه حدود ۶۱ و ۵۹ درصد جت‌ها (به ترتیب \bbbb و \ttbb) $\Delta R_{b\bar b}<0.3$ دارند، ولی در سیگنال فقط حدود ۱۲ و ۱۳ درصد. سپس جت‌ها را بر اساس $\Delta R_{b\bar b}$ دسته‌بندی کردیم و میانگین $\Delta R$ تا محور جت را برای ده ذره پرانرژی، با وزن $p_T$ ذره، حساب کردیم (جدول~\ref{tab:drbb}). در دو دسته $0.3\le\Delta R_{b\bar b}<0.5$ سیگنال و پس‌زمینه تقریباً یکی‌اند (مثلاً \lr{0.139} در برابر \lr{0.133} برای \Zjet و \bbbb در دسته اول)، و بیشتر پس‌زمینه در دسته $\Delta R_{b\bar b}<0.3$ است که میانگین آن فقط حدود \lr{0.08} تا \lr{0.11} است. پس بخش بزرگی از تفاوت قله $\Delta R$ ذره‌ها با تفاوت $\Delta R_{b\bar b}$ توضیح داده می‌شود. \textcolor{red}{در دسته‌های $\Delta R_{b\bar b}<0.3$ و بالای \lr{0.5} تفاوت باقی می‌ماند و هنوز توجیهی برای آن نداریم.}

\begin{table}[H]
\centering
\caption{میانگین $\Delta R$ تا محور جت برای ده ذره پرانرژی (وزن‌دار با $p_T$) در دسته‌های $\Delta R_{b\bar b}$؛ داخل پرانتز تعداد جت. جت‌ها با هر دو $b$ در جت و $p_T>400$ و $|\eta|<2$ و $n_\mathrm{const}\ge3$.}
\label{tab:drbb}
\begin{tabular}{lcccc}
\hline
دسته $\Delta R_{b\bar b}$ & \Zjet & \bbbb & \ttZ & \ttbb \\
\hline
$0$ تا $0.3$ & $0.106\ (282)$ & $0.083\ (2799)$ & $0.115\ (373)$ & $0.085\ (1506)$ \\
$0.3$ تا $0.4$ & $0.139\ (686)$ & $0.133\ (557)$ & $0.150\ (1062)$ & $0.139\ (334)$ \\
$0.4$ تا $0.5$ & $0.141\ (714)$ & $0.148\ (447)$ & $0.162\ (863)$ & $0.156\ (234)$ \\
$0.5$ تا $0.6$ & $0.104\ (323)$ & $0.166\ (335)$ & $0.118\ (288)$ & $0.182\ (219)$ \\
$0.6$ تا $0.8$ & $0.081\ (328)$ & $0.198\ (423)$ & $0.090\ (210)$ & $0.196\ (233)$ \\
\hline
\end{tabular}
\end{table}

\widefig{cons_10.pdf}{$\eta$ چرخیده برای ۱۰ ذره پرانرژی.}{fig:cons-etarot}
\widefig{cons_13.pdf}{$\varphi$ چرخیده برای ۱۰ ذره پرانرژی.}{fig:cons-phirot}
در قاب چرخیده ذره اول در مبدأ ثابت است و ذره دوم روی محور $\varphi$ قرار می‌گیرد؛ به همین دلیل $\eta$ چرخیده ذره دوم در حد خطای عددی (از مرتبه $10^{-16}$) صفر است و $\varphi$ چرخیده ذره اول ثابت است (شکل‌های~\ref{fig:cons-etarot} و~\ref{fig:cons-phirot}). این دو را نباید به‌عنوان ویژگی آموزنده شمرد. برای ذره دوم $\varphi$ چرخیده همه‌جا نامثبت است و قله‌ای تیز در صفر دارد؛ این علامت از تعریف کد می‌آید: ذره دوم روی محور منفی $\varphi$ چرخیده قرار داده می‌شود. برای ذرات سوم تا دهم، هر دو توزیع قله‌ای تیز در صفر و دنباله‌هایی دارند که نسبت سیگنال به پس‌زمینه آن‌ها دو برآمدگی نشان می‌دهد: در $\varphi_\mathrm{rot}$ در حدود $-0.35$ (ارتفاع حدود ۲ تا ۳) و $+0.3$ (حدود ۲) و در $\eta_\mathrm{rot}$ در حدود $\pm0.3$ تا $\pm0.5$ (حدود ۲)، با فرورفتگی کم در صفر (حدود \lr{0.8} تا ۱). یعنی در سیگنال، ذره‌های بیشتری در فاصله حدود \lr{0.3} تا \lr{0.5} از ذره اول، خارج از قله‌ای که ذره‌های همسایه آن در آن جمع‌اند، دیده می‌شوند. این رفتار با ساختار دو زیرجت سازگار است، چون فاصله دو هسته در بخش بالا حدود \lr{0.46} برای $p_T=400$ (و کمتر برای $p_T$ بزرگ‌تر) برآورد شد. \textcolor{red}{این تفسیر فرضیه است و با ربط دادن ذره‌ها به کوارک‌های $b$ سطح پارتون آزمون نشده است؛ معنی فیزیکی برآمدگی‌های $\eta_\mathrm{rot}$ را هم هنوز نمی‌دانیم.} دو جفت در این ویژگی‌ها هم‌جهت‌اند.

\widefig{cons_15.pdf}{$\cos\theta$ ده ذره پرانرژی.}{fig:cons-cos}
\widefig{cons_16.pdf}{$\cos\theta_\mathrm{rel}$ (زاویه نسبت به جت) برای ده ذره پرانرژی.}{fig:cons-cosrel}
شکل~\ref{fig:cons-cos} \lr{U}‌شکل است. اگر $\cos\theta$ زاویه قطبی ذره نسبت به محور پرتو باشد، آنگاه $\cos\theta=p_z/|\vec p|=\tanh\eta$ است و چگالی آن $\rho(\cos\theta)=\rho(\eta)\,/(1-\cos^2\theta)$ می‌شود که در لبه‌ها بالا می‌رود؛ این با \lr{U}‌شکل‌بودن سازگار است. رفتار نسبت‌ها دقیقاً مثل $\eta$ است: نسبت \Zjet به \bbbb از حدود \lr{0.8} تا \lr{0.9} در مرکز تا حدود \lr{1.3} در لبه، و نسبت \ttZ به \ttbb از حدود \lr{1.1} در مرکز تا حدود \lr{0.7} در لبه. پس این ویژگی هم از همان تفاوت $\eta$ جت‌ها می‌آید و دو جفت خلاف جهت‌اند. $\cos\theta_\mathrm{rel}$ (شکل~\ref{fig:cons-cosrel}) کسینوس زاویه ذره با محور جت است و همه‌جا نزدیک ۱ جمع شده است. برای زاویه‌های کوچک $\cos\theta_\mathrm{rel}\approx1-\theta^2/2$ است، پس $\theta\approx\sqrt{2(1-\cos\theta_\mathrm{rel})}$. مقدار $\cos\theta_\mathrm{rel}\approx0.98$ تا \lr{0.995} یعنی $\theta\approx0.1$ تا \lr{0.2}، که همان بازه برآمدگی $\Delta R$ در شکل~\ref{fig:cons-dr} است، و مقدار حدود \lr{0.92} تا \lr{0.95} یعنی $\theta\approx0.3$ تا \lr{0.4} که در آن نسبت زیر ۱ است. نسبت‌ها هم همین‌طورند: زیر ۱ برای $\cos\theta_\mathrm{rel}<0.96$، تا حدود \lr{1.8} (\ttZ/\ttbb) و حدود \lr{1.3} (\Zjet/\bbbb) نزدیک \lr{0.99}. پس این ویژگی همان اطلاعات $\Delta R$ را با مقیاس دیگری می‌دهد.

\subsection{نوع ذره}
\widefig{cons_17.pdf}{نوع ذره ($|\mathrm{pdgId}|$) برای ۱۰ ذره پرانرژی.}{fig:cons-pdg}
ترکیب نوع ذره‌ها (الکترون ۱۱، میون ۱۳، فوتون ۲۲، هادرون خنثی ۱۳۰، پیون باردار ۲۱۱، کائون ۳۲۱، پروتون ۲۲۱۲) در چهار نمونه تقریباً یکسان است (شکل~\ref{fig:cons-pdg}). برای ذره اول، پیون باردار حدود \lr{0.27} و فوتون و هادرون خنثی هر کدام حدود \lr{0.2} تا \lr{0.25} است، و برای ذرات بعدی سهم پیون باردار به حدود \lr{0.4} تا \lr{0.45} بالا می‌رود. نسبت‌ها در حدود ۱ می‌مانند و تنها برای ذره اول در دسته الکترون و میون و فوتون تفاوتی حدود ۱۰ تا ۳۰ درصد دیده می‌شود. چون دسته الکترون و میون سهم کوچک (حدود \lr{0.05}) دارد و خطای آماری محاسبه نشده است، نمی‌توانیم این تفاوت را معنی‌دار بدانیم. پس این ویژگی اطلاعات اضافه‌ای برای جداسازی نمی‌دهد.

\subsection{جمع‌بندی}
ویژگی‌ها را با دو معیار می‌توان دسته‌بندی کرد: اندازه انحراف نسبت از ۱، و اینکه دو جفت هم‌جهت‌اند یا خلاف جهت. ویژگی‌هایی که فقط از $\eta$ جت می‌آیند ($p_z$، $E$، $\eta$ و $\cos\theta$) انحراف‌های بزرگی دارند ولی دو جفت در آن‌ها خلاف جهت‌اند؛ این همان خطر میان‌بر بخش~\ref{sec:set4} است، چون شناساگر می‌تواند به‌جای ساختار \Zbos، توزیع $\eta$ را یاد بگیرد. ویژگی‌های نسبی هندسی ($\eta-\eta_\mathrm{jet}$، $\varphi-\varphi_\mathrm{jet}$، $\Delta R$، $\cos\theta_\mathrm{rel}$ و مختصات چرخیده ذرات سوم تا دهم) در هر دو جفت یک جهت دارند و با ساختار دو‌هسته‌ای سازگارند؛ این‌ها نامزدهای اصلی ورودی شناساگر هستند. $p_T/p_{T,\mathrm{jet}}$ و $E/E_\mathrm{jet}$ انحراف متوسطی دارند و در دو جفت هم‌جهت‌اند، ولی به‌صورت کامل آزموده نشده‌اند. $p_x$ و $p_y$ و $\varphi$ و نوع ذره اطلاعات اضافه‌ای ندارند. $\eta$ چرخیده ذره دوم و $\varphi$ چرخیده ذره اول با ساخت ثابت‌اند.
